\clearpage
\thispagestyle{empty}
\SourceLeaf{417}\pdfbookmark[0]{Perturbations of Pallas}{pallas}
\vspace*{1.6in}
\begin{center}{\Large PERTURBATIONS OF PALLAS.\par}\end{center}
\clearpage
\fancyhead[L]{\small PERTURBATIONS OF PALLAS}
\SourceLeaf{419}\pdfbookmark[1]{I. Correspondence}{pallas-briefwechsel}
\begin{center}
{\large [I. CORRESPONDENCE.]\par}
\bigskip
\pdfbookmark[2]{Gauss to Olbers, 3 January 1806}{pallas-brief-18060103}
\textsc{Gauss} to \textsc{Olbers}, Braunschweig, 3 January 1806.
\end{center}
\noindent\hspace*{\parindent}.\quad.\quad.\quad.\quad.\quad Some chance circumstances led me a while ago to begin working out various investigations on the theory of interpolation\HistoricalNoteEight{See vol. III, p. 265.]}. I thought I could fit this into one or a couple of sheets. It has now grown to 4, which might fill 8 in print. It is precisely the latter, larger half that I find essential and very useful in my method for calculating perturbations. I would therefore like an opportunity to have that paper printed somewhere beforehand.\quad.\quad.\quad.\quad.\quad.
\par\bigskip
\begin{center}
\pdfbookmark[2]{Gauss to Olbers, 6 August 1810}{pallas-brief-18100806}
\textsc{Gauss} to \textsc{Olbers}, Göttingen, 6 August 1810.
\end{center}
\noindent\hspace*{\parindent}.\quad.\quad.\quad.\quad.\quad I have begun writing a paper\HistoricalNoteDouble{See vol. VI, p. 1.]} for the Society on my calculations concerning Pallas. Besides the numerical results, it will contain my method for determining a planetary orbit from 4 oppositions (more precisely, the elliptic elements from 4 longitudes in the orbit), as well as the shortcuts I use in the method of least squares. These may be of some interest to you. If the deadline for the Paris prize problem were 4 months further away, I might still enter now. As things stand, the time is too short, and for 10 months I have had neither the inclination nor the courage for any \SourceLeaf{420}somewhat lengthy, sustained work. Even so, I now intend, in some 4 or 5 separate memoirs for the Society, to work out all parts of the method by which I believe the perturbations of Pallas can be calculated efficiently.
\par\bigskip
\begin{center}
\pdfbookmark[2]{Gauss to Olbers, 24 October 1810}{pallas-brief-18101024}
\textsc{Gauss} to \textsc{Olbers}, 24 October 1810.
\end{center}
\noindent\hspace*{\parindent}.\quad.\quad.\quad.\quad.\quad Soon I shall read to the Society the paper\HistoricalNoteEight{See vol. VI, p. 1.]} on Pallas that I believe I once mentioned to you. You will then find a detailed account of it in our Anzeigen. It also contains some theory. I develop my method, already in use for several years, for determining an elliptical orbit from 4 oppositions, a kind of supplement to the 2nd section of the 2nd book of my Theoria. I also explain a practical device that considerably shortens the otherwise laborious elimination in the method of least squares.\quad.\quad.\quad.\quad.\quad.\quad I may therefore soon start calculating the perturbations of one of the new planets more fully. If I knew that the Institute would extend the deadline for the prize problem once more, I should be inclined to choose {\AstroSymbols ⚴}. Otherwise {\AstroSymbols ⚳} must rightly take precedence. Its 8th {\AstroSymbols ☍} will soon be observed, so that a more satisfactory result can be expected there too. With {\AstroSymbols ⚴} I should probably begin by treating the elements as variable and determining their perturbations over 7 years by quadratures. For this I have devised formulas that seem to me somewhat more convenient than \textsc{Laplace}'s. If I then succeed, as I expect to, in reconciling the 6 oppositions observed so far, which can no longer be fitted into a single ellipse at all, I should use the method I devised 4 or 5 years ago to calculate the perturbations in the usual form, as periodic perturbations of the longitude, latitude and radius vector, or better still of its logarithm.\quad.\quad.\quad.\quad.\quad.
\par
